\begin{abstract}
Machine unlearning promises to remove specific knowledge from a trained language model without the cost of
retraining, a capability demanded by privacy regulation, copyright, and AI safety. Yet whether an unlearned
model has genuinely \emph{forgotten} a fact, rather than merely learned to \emph{hide} it, is
surprisingly hard to verify. Unlearning is evaluated almost entirely at the output, with lexical and
behavioural metrics such as accuracy, ROUGE, gold-answer probability, and membership inference; a growing
body of work shows these metrics are fragile, as quantisation, a few steps of relearning, or a direct probe
of the model's internal states can each revive the ``forgotten'' knowledge. The gap is between
\emph{suppression} at the output and \emph{erasure} inside the model.

This thesis closes that gap on the evaluation side by reading the model's internal computation with
mechanistic interpretability. Building on the Hubble benchmark, whose controlled fact insertions provide a
never-knew and a knows reference for every fact, we integrate a corrected, boundary-safe logit lens that
records the rank of the correct answer at each layer, and from this trajectory we derive a dual-anchor
distance, \(D_\text{mech}\), that measures whether a fact's internal trajectory looks more like the knows
reference or the never-knew reference, independently of what the model outputs. Under the field's own
faithfulness meta-evaluation based on OpenUnlearning, on the Hubble-8B suite \(D_\text{mech}\) is the only
metric, mechanistic or lexical, that is at once strongly knowledge-discriminative and substantially
paraphrase-stable, where every lexical metric that discriminates well collapses under a meaning-preserving
rewording while \(D_\text{mech}\) shifts far less. Its composite faithfulness score is about eighteen
times the best lexical metric, an ordering that we show is robust to the choice of aggregation formula. We
then replicate the meta-evaluation on a second benchmark and model family (OpenUnlearning TOFU, at 1B, 3B,
and 8B), where \(D_\text{mech}\) is again faithful and, at every scale, the best internal metric and better
than every generation-overlap and answer-probability metric, though a few likelihood-based aggregates
become competitive there. We further characterise all three interpretability lenses on the same facts,
explaining empirically why the calibrated tuned lens ``skips ahead'' and cannot ground the metric, and we
report a scale-dependent verdict on the four unlearning methods that is invisible at the output. Read
together, the results indicate that the ``forgotten'' facts studied here remain internally accessible,
suppressed at the output rather than erased from the representation as far as the applied readouts can
determine, and that internal readouts yield a more faithful way to tell the difference.
\end{abstract}
